%\documentclass[10pt,respuestas,a4]{examenEPN}
\documentclass[10pt,a4]{examenEPN0}
\usepackage[T1]{fontenc}
\usepackage{multicol}
\usepackage{blindtext}

%\documentclass[10pt,a4]{examenEPN0}
\examen{Prueba 01}

\usepackage{comandosEPN}

\usepackage{empheq}

\newcommand{\itemcp}{\stepcounter{enumi}\item[ \textcolor{red}{\theenumi.$^*$}] }

\newcommand{\iitemcp}{\stepcounter{enumii}\item[ \textcolor{red}{\theenumii)$^*$}] }


\materia{Álgebra Lineal}
%\autor{Departamento de Formación Básica}
\autor{GR \hspace{1cm}.}
%\fecha{Semestre 2026-A}
\fecha{Nombre: \hspace{9cm} Cédula:}
\logouno[1.6cm]{Logo_EPN.eps}
\logodos[2.6cm]{Logo_DFB.eps}

\newcommand{\la}{\lambda}

%Datos
\newcommand{\datosAlgebra}[2]{
\EnExamen
{
\vspace{5mm}
    \noindent
    Nombres: \hrulefill\
    Carrera: \rule{4.5cm}{0.4pt}\
    \ifthenelse{\equal{#1}{}}%
    {Paralelo: \underline{\parbox{1.8cm}{\parbox{1cm}{}}}\\[3mm]}
    {Paralelo: \underline{\parbox{1.8cm}{\centering #1}}\\[3mm]}
    C\'edula: \rule{5cm}{0.4pt}\
    Firma: \hrulefill\
    Fecha: \underline{\parbox{2.1cm}{\centering #2}}
}}


% \newcommand{\indicacionesPrueba}{%
% \EnExamen{
% \begin{indicaciones}
%     \item La prueba tiene una duraci\'on  de \textbf{120 minutos.}
%     \item No se permite la utilizaci\'on de dispositivos electr\'onicos (celular, tablets, calculadora, etc.)
%     \item Leer detenidamente los enunciados, \textbf{justificar la respuesta para acceder al puntaje total}.
%     \item La prueba es completamente individual. Los estudiantes deberán dejar sus teléfonos celulares guardados y no podrán tenerlos consigo (por ejemplo, en los bolsillos de su ropa) durante toda la evaluación. En caso de detectarse algún acto de deshonestidad académica o el incumplimiento de esta norma, no podrán seguir rindiendo la prueba y su calificación será \textbf{cero}.
% \end{indicaciones}
% \pagestyle{empty}
% }}

\newcommand{\indicacionesPrueba}{%
\EnExamen{
\begin{indicaciones}
    \item La prueba tiene una duraci\'on  de \textbf{60 minutos.}
    \item Leer detenidamente los enunciados, \textbf{justificar la respuesta} para acceder al puntaje total.
    \item No se permite la utilizaci\'on de dispositivos electr\'onicos (celular, tablets, calculadora, etc.)
    \item La prueba es completamente individual
    \item No se permite portar teléfonos celulares, relojes inteligentes ni otros dispositivos electrónicos durante el examen, aun cuando permanezcan en los bolsillos. Estos dispositivos deberán guardarse en la mochila antes del inicio de la evaluación y esta deberá mantenerse cerrada durante todo el examen.
    \item He leído las indicaciones y acepto que el incumplimiento de estas se penaliza con \textbf{cero}. Firma:...................................

    
\end{indicaciones}
\pagestyle{empty}
}}

\newcommand{\indicacionesExamen}{
\EnExamen{\begin{indicaciones}
    \item La evaluación tiene una duraci\'on  de \textbf{120 minutos.}
    \item No se permite la utilizaci\'on de dispositivos electr\'onicos (celular, tablets, calculadora, etc.)
    %\item Leer detenidamente los enunciados, justificar la respuesta para acceder al puntaje total.
    \item El examen es completamente individual. En caso de detectarse alg\'un acto de deshonestidad acad\'emica no podr\'a seguir rindiendo la prueba y su calificaci\'on ser\'a \textbf{cero}.
\end{indicaciones}
}\pagestyle{empty}
}

%% Imprimir las una hoja A4, dos veces en A5
\newcommand{\imprimirDoble}
{%
\usepackage{pgfmorepages}
\pgfmorepagesloadextralayouts
\pgfpagesuselayout{repeated 2-up}[a4paper,landscape]
\thispagestyle{empty}
}


%Este sí se puede comentar
\newenvironment{solucion}{\color{blue}\proof[Solución]}{\endproof\color{black}}




% \renewcommand{\@autor}{
% ---\\
% \phantom{.} 
% ...}
% \makeatother


%\imprimirDoble

\begin{document}
\encabezado
\datosAlgebra{GR5M}{21/10/2026}

%Poner el grupo y la fecha, o dejar las llaves vacías.
\indicacionesPrueba

\begin{preguntas}

% PREGUNTA 1

\item Considere el sistema homogéneo $Ax=\mathbf{0}$, donde
\[
    A = \begin{pmatrix} 1 & 0 \\ 1 & 1 \\ 0 & 1 \end{pmatrix},
    \qquad x = \begin{pmatrix} x_1 \\ x_2 \end{pmatrix}.
\]
Calcule el rango de $A$ y compárelo con el número de incógnitas para determinar si el sistema tiene solución única o infinitas soluciones. Justifique su respuesta y escriba el conjunto solución.
    


\puntaje{5}
% Espacio de respuesta hasta el final de la página.
\par\vfill\newpage


\item Un estudiante quiere demostrar o refutar que $(AB)^2 = A^2 B^2$ para toda pareja de matrices cuadradas $A$ y $B$ del mismo orden. Para esto escribe el siguiente razonamiento:
\par\medskip
\noindent\colorbox{gray!15}{%
    \begin{minipage}{\dimexpr\linewidth-2\fboxsep\relax}
        \normalfont\normalsize\raggedright
        \textbf{Demostración propuesta por el estudiante}
        \begin{enumerate}[label=\arabic*.,leftmargin=*,topsep=4pt,itemsep=2pt,parsep=0pt]
            \item $(AB)^2 = (AB)(AB)$.
            \item $(AB)(AB)=A(BA)B$.
            \item $A(BA)B= A(AB)B$.
            \item $A(AB)B = A^2 B^2$.
        \end{enumerate}
        \textbf{Conclusión del estudiante:} Por lo tanto, la propiedad
        $(AB)^2 = A^2 B^2$ siempre es verdadera para matrices cuadradas
        del mismo orden.
    \end{minipage}%
}
\par\medskip


    
    Determine si la demostración es correcta. Si lo es, justifique cada paso. Si no lo es, identifique el paso incorrecto y presente un ejemplo con matrices de orden $2\times2$ que refute la afirmación inicial.


\puntaje{5}
% Espacio de respuesta hasta el final de la página.
\par\vfill\newpage

\item 

Considere el siguiente sistema de ecuaciones, donde $k \in \mathbb{R}$ es un parámetro:
\[
\begin{cases}
x_1 + x_2 + x_3 = 1 \\
x_2 + kx_3 = 2 \\
x_1 + 2x_2 + 2x_3 = 3
\end{cases}
\]
\begin{enumerate}
    \item Escriba la matriz ampliada $(A \mid b)$ asociada al sistema y aplique, en ese orden, las operaciones elementales por fila $F_3 \leftarrow F_3-F_1$ y $F_3 \leftarrow F_3-F_2$. En la segunda operación, utilice la fila $F_3$ obtenida en la primera. \puntaje{1}

    \item Determine los rangos de $A$ y de $(A \mid b)$ según el valor de $k$. Aplique el teorema de Rouché--Frobenius para indicar, en cada caso, si el sistema tiene solución única, infinitas soluciones o ninguna solución. \puntaje{4}
\end{enumerate}


% Espacio de respuesta hasta el final de la página.
\par\vfill\newpage


\item En un modelo simplificado para una obra, se necesitan $100\,\text{m}^3$ de una mezcla de cemento, arena y grava. Suponga que los volúmenes se suman al combinar los materiales. La mezcla final debe contener exactamente $28\,\text{m}^3$ de cemento, $42\,\text{m}^3$ de arena y $30\,\text{m}^3$ de grava.

Se dispone de tres mezclas ya preparadas, llamadas $M_1$, $M_2$ y $M_3$, con la siguiente composición en volumen:
\begin{itemize}
    \item Mezcla $M_1$: $30\%$ de cemento, $40\%$ de arena y $30\%$ de grava.
    \item Mezcla $M_2$: $20\%$ de cemento, $50\%$ de arena y $30\%$ de grava.
    \item Mezcla $M_3$: $40\%$ de cemento, $30\%$ de arena y $30\%$ de grava.
\end{itemize}
Se desea establecer cuántos metros cúbicos de cada mezcla preparada se deben utilizar para obtener la mezcla final con las cantidades requeridas de cemento, arena y grava.

\begin{enumerate}
    \item Defina las tres incógnitas e indique qué representa cada una y en qué unidad se mide. Considere que las cantidades utilizadas deben ser no negativas. \puntaje{1}
    \item Plantee una ecuación para la cantidad de cemento, otra para la cantidad de arena y otra para la cantidad de grava. \puntaje{3}
    \item Escriba el sistema en la forma matricial $Ax=b$ e identifique la matriz $A$ y los vectores $x$ y $b$. \puntaje{1}
\end{enumerate}
No es necesario resolver el sistema.



% El espacio restante de la última página queda disponible para responder.






\vfill

\end{preguntas}



\final
\end{document}
